\section{官方幻灯片证据链：从 Agent Demo 到生产产品}

Clay Bavor 的课堂主线可以压缩成一句话：生产 Agent 的难度主要藏在模型之外。用户只看到一次自然语言交互，团队却要同时解决业务状态、工具执行、记忆、渠道连续性、语音延迟、测试仿真、可观测性与安全策略。本节按 40 页官方 deck 的顺序保留完整证据链，再由后续章节把这些经验转成架构和运营手册。

\subsection{市场速度、体验迁移与单 Agent 世界}

slides 1--8 从 Sierra 的业务背景出发，把 AI 类比为五倍速播放的互联网浪潮。互联网把体验移到网站，移动互联网把体验移到 App，Agent 则把自然语言对话变成界面。结果不是再增加一个客服渠道，而是电话、邮件和聊天逐渐坍缩为同一个有连续记忆的 Agent；商业模式也从按席位或调用次数收费，转向为完成的工作付费。

\lecturefigure{slides-images/slide-001.png}{官方 slide 1：Practical Lessons from Deploying Real-World AI Agents}
\lecturefigure{slides-images/slide-002.png}{官方 slide 2：Sierra 的业务、团队与客户背景}
\lecturefigure{slides-images/slide-003.png}{官方 slide 3：真实客户与 Agent 生态概览}
\lecturefigure{slides-images/slide-004.png}{官方 slide 4：AI adoption 如同五倍速播放的互联网曲线}
\lecturefigure{slides-images/slide-005.png}{官方 slide 5：Internet、Mobile 与 AI 分别创造新体验}
\lecturefigure{slides-images/slide-006.png}{官方 slide 6：从 multi-channel 转向 single-agent world}
\lecturefigure{slides-images/slide-007.png}{官方 slide 7：Conversation is the interface}
\lecturefigure{slides-images/slide-008.png}{官方 slide 8：为完成的工作付费}

\begin{knowledgebox}{读图：渠道坍缩不是把所有入口接到同一个聊天框}
真正统一的是客户身份、长期状态和行动权限。若电话 Agent 不知道邮件承诺，或聊天 Agent 不能继续上次退换货流程，界面虽然统一，体验仍是碎片化的。outcome pricing 又要求系统对任务完成负责，因此“说得像完成了”必须与业务系统中的可验证状态分开。
\end{knowledgebox}

\begin{importantbox}{课堂提示：先定义完成，再谈定价}
按结果收费会把供应商激励与客户价值对齐，但前提是结果可判定、责任边界清楚、外部依赖可记录。退款到账、设备激活和预约确认可由系统状态验证；“关系更好”或“客户满意”则需要长期指标，不能用一次对话自报替代。
\end{importantbox}

\subsection{Agent Iceberg：模型只是水面以上}

slides 9--13 用冰山解释 production gap。水面上的 LLM、RAG 和 tool use 容易做出 demo；水面下还包括复杂工作流、审计、可观测性、控制、数据平台、评测、部署和运营。把 Agent 从 technology 变成 product，需要像 1997 年的网站开发一样逐步沉淀工具链，同时坚持“simple, not simplistic”：用户路径简单，内部系统不能假装现实也简单。

\lecturefigure{slides-images/slide-009.png}{官方 slide 9：Agent Iceberg 的可见层}
\lecturefigure{slides-images/slide-010.png}{官方 slide 10：Agent Iceberg 的生产系统全貌}
\lecturefigure{slides-images/slide-011.png}{官方 slide 11：Agent 开发仍处于类似 1997 年的早期阶段}
\lecturefigure{slides-images/slide-012.png}{官方 slide 12：从 agents as technology 转向 agents as product}
\lecturefigure{slides-images/slide-013.png}{官方 slide 13：Simple, not simplistic}

\begin{knowledgebox}{术语消化：冰山下的四类能力}
\begin{tabular}{p{2.8cm}p{4.2cm}p{5.2cm}}
\toprule
\textbf{能力} & \textbf{解决的问题} & \textbf{生产含义} \\
\midrule
Orchestration & 多步任务如何推进 & 显式状态机、回退、重试与人工升级 \\
Observability & 为什么成功或失败 & trace、指标、版本和逐步决策证据 \\
Controls & Agent 可以做什么 & 权限、策略、审批与业务规则 \\
Operations & 系统如何持续改进 & 仿真回归、线上告警、知识更新与事故复盘 \\
\bottomrule
\end{tabular}
\end{knowledgebox}

\subsection{从事务处理到长期关系}

slides 14--16 对比两种产品目标。事务型 Agent 完成退货、故障排查、套餐升级或账号恢复；关系型 Agent 还要跨渠道持续参与、记住互动、接入企业数据并主动服务。后者要求 memory 不只是聊天摘要，而是带来源、时效和权限的客户状态；proactive engagement 也必须受同意、频率和合规约束。

\lecturefigure{slides-images/slide-014.png}{官方 slide 14：今天的 Agent 多数仍是事务型}
\lecturefigure{slides-images/slide-015.png}{官方 slide 15：最好的 Agent 不只解决 case，而是建立关系}
\lecturefigure{slides-images/slide-016.png}{官方 slide 16：跨渠道记忆、企业数据与主动服务}

\begin{warningbox}{长期记忆的错误会比单轮幻觉更危险}
一次错误回答可能被用户纠正；错误客户状态若写入长期记忆，会污染后续所有渠道和主动行动。生产系统应区分事实、推断与偏好，记录来源和更新时间，并允许客户查看、更正和删除。warm start 的目标是减少重复说明，不是让旧判断永久支配新情境。
\end{warningbox}

\subsection{Agent Development Life Cycle 与 Build-with}

slides 17--22 提出 build-test-optimize 循环，并把 build-or-buy 改写为 build-with：企业需要保留业务规则和体验所有权，同时复用成熟平台的底层能力。每个渠道已经数字化，因此同一逻辑可部署到电话、邮件和聊天；但真正的“一次构建”依赖共享 Agent Data Platform，把 memory、客户数据、智能决策和主动互动放在统一状态层。

\lecturefigure{slides-images/slide-017.png}{官方 slide 17：Agent Development Life Cycle}
\lecturefigure{slides-images/slide-018.png}{官方 slide 18：Build with，而非简单 build or buy}
\lecturefigure{slides-images/slide-019.png}{官方 slide 19：Every channel is digital}
\lecturefigure{slides-images/slide-020.png}{官方 slide 20：Build once, deploy everywhere}
\lecturefigure{slides-images/slide-021.png}{官方 slide 21：Technology 到 product 的体验差距}
\lecturefigure{slides-images/slide-022.png}{官方 slide 22：Agent Data Platform 的四个组成}

\begin{knowledgebox}{读图：Build-test-optimize 是闭环而非三个部门}
构建阶段必须产出可测状态和 trace；测试阶段的失败要自动回流为仿真场景；优化阶段改变 prompt、知识、策略或工具后，必须重新跑同一回归集。若三个阶段使用不同数据格式和负责人手工交接，改进速度会被治理成本吞掉。
\end{knowledgebox}

\subsection{Voice：低时延只是第一层门槛}

slides 23--26 拆解 voice pipeline：speech-to-text 转写，LLM 和工具完成 reasoning，text-to-speech 合成响应。真实质量还受口音、噪声、打断、停顿、实体读法和情绪影响。slide 24 直接指出 WER 对 voice agent 很糟，因为普通拼写错误与电话号码、网址或订单号错误的业务代价完全不同；slide 25 因此要求 locale-aware entity synthesis 和自然节奏。

\lecturefigure{slides-images/slide-023.png}{官方 slide 23：Voice Agent 的 STT、reasoning 与 TTS 架构}
\lecturefigure{slides-images/slide-024.png}{官方 slide 24：WER 无法充分度量业务语音质量}
\lecturefigure{slides-images/slide-025.png}{官方 slide 25：自然实体合成、韵律、打断与情绪}
\lecturefigure{slides-images/slide-026.png}{官方 slide 26：什么才算自然可用的声音}

\begin{importantbox}{Voice 质量函数}
可把用户体验粗略写成
\[
Q_{\text{voice}}
=w_1Q_{\text{intent}}+w_2Q_{\text{entity}}+w_3Q_{\text{timing}}+w_4Q_{\text{naturalness}}-w_5E_{\text{action}},
\]
其中 intent 是意图理解，entity 是关键实体准确率，timing 是首包延迟与打断响应，naturalness 是韵律和音色，$E_{\text{action}}$ 是错误业务行动的代价。权重 $w_i$ 必须按场景设定，不能用一个通用 WER 替代。
\end{importantbox}

\subsection{Test：从用户模拟器到 tau-bench}

slides 27--33 把测试视为新软件范式的核心。早期用户模拟器只能生成对话，真实 Agent 还要面对工具、用户状态、业务政策、数据库和长轨迹。tau-bench 把 realistic domain、user simulator 和 objective evaluation 结合，并通过 pass@k 检查同一任务多次运行的一致性；社区采用说明可复现 benchmark 能成为产业共同语言。Sims 与 voice sims 则把生产失败转为可重复场景。

\lecturefigure{slides-images/slide-027.png}{官方 slide 27：New testing for new software}
\lecturefigure{slides-images/slide-028.png}{官方 slide 28：2023 年最早期用户与 Agent 模拟器}
\lecturefigure{slides-images/slide-029.png}{官方 slide 29：真实 Agent 测试需覆盖工具、政策、状态和规模}
\lecturefigure{slides-images/slide-030.png}{官方 slide 30：tau-bench 的领域、模拟器、目标评测与一致性}
\lecturefigure{slides-images/slide-031.png}{官方 slide 31：tau-bench 与 tau2-bench 的社区采用}
\lecturefigure{slides-images/slide-032.png}{官方 slide 32：生产对话场景 Sims}
\lecturefigure{slides-images/slide-033.png}{官方 slide 33：Voice Sims compilation}

\begin{knowledgebox}{读图：pass@k 在 Agent 场景衡量什么}
若一个任务独立运行 $k$ 次，pass@k 可描述至少一次成功或全部成功等不同口径。生产可靠性更关心 repeated success：同一退款流程不能十次中只成功一次。报告时必须说明是“至少一次成功”、平均成功率，还是 $k$ 次全部通过，并同时记录失败类型和状态污染。
\end{knowledgebox}

\subsection{Optimize：Traces、Insights 与 Expert Answers}

slides 34--36 展示从运行数据反哺系统。trace 把一次会话中的模型决策、工具调用和业务状态串起来；Insights Explorer 从大量对话中聚类原因与机会；Expert Answers 让领域专家把正确处理方式沉淀为可复用知识。这里的“用 AI 改进 AI”不是让模型无监督改写自己，而是提高问题发现和知识运营的吞吐量。

\lecturefigure{slides-images/slide-034.png}{官方 slide 34：生产 Agent 的可观测 Traces}
\lecturefigure{slides-images/slide-035.png}{官方 slide 35：Optimize -- Use AI to improve AI}
\lecturefigure{slides-images/slide-036.png}{官方 slide 36：Insights Explorer 与 Expert Answers}

\begin{warningbox}{优化闭环必须防止自证循环}
如果同一个模型生成回答、评判回答并改写规则，偏差可能被放大。高影响问题应由真实业务结果、人工专家或独立 evaluator 校准；Insights 给出调查线索，不自动等于根因。每次知识更新还应绑定来源、批准者、版本和回归结果。
\end{warningbox}

\subsection{Layered Guardrails：安全是系统拓扑}

slides 37--40 用攻击案例和分层图结束。提示注入可以利用外部内容操纵 Agent，因此防护不能只放在输入过滤器。业务系统、读写权限、规则、政策、supervisor、检索、知识与记忆都要有独立边界；输入和输出还需做内容与行动检查。最后的招聘页虽是行政收尾，也表明这套工程仍处于快速建设期。

\lecturefigure{slides-images/slide-037.png}{官方 slide 37：安全章节转场}
\lecturefigure{slides-images/slide-038.png}{官方 slide 38：通过外部内容攻击 Agent 的现实案例}
\lecturefigure{slides-images/slide-039.png}{官方 slide 39：Business Systems、Guardrails、Policies、Supervisors、Retrieval、Knowledge 与 Memory}
\lecturefigure{slides-images/slide-040.png}{官方 slide 40：课程收尾}

\begin{knowledgebox}{读图：每一层都应能独立拒绝危险动作}
输入层识别恶意指令，检索层过滤不可信来源，策略层判断任务是否允许，工具层实施最小权限，supervisor 对高风险动作审批，输出层检查泄露和误导，业务系统本身仍需事务约束。任何一层失守都不应直接导致不可逆行动，这就是 defense in depth。
\end{knowledgebox}

\begin{warningbox}{课堂提醒：Guardrail 不是替模型擦屁股}
安全层无法弥补没有权限模型、没有审计日志或没有业务幂等性的架构。真实部署要假设模型会被诱导、会误解、会超时；系统必须把失败限制在可恢复范围，并让人能够解释、撤销和修复。
\end{warningbox}

\subsection{本章小结}
40 页 deck 给出了一套完整生产方法：把 conversation 视为统一界面，以 outcome 定义价值；认识 Agent Iceberg，把 memory、data、controls 和 operations 纳入产品；以 build-test-optimize 迭代，针对 voice 建专用指标，用仿真与 trace 建可靠性闭环，最后用多层 guardrails 限制失败半径。后续章节将这些证据转为实施细节。

